\documentclass[11pt]{article}
\usepackage[a4paper,left=5.57cm,right=1.8cm,top=4.58cm,bottom=3.5cm,
  headheight=14pt,headsep=18pt]{geometry}
% Margins match the reference photo, measured on its de-warped image: the
% text block starts 158 pt from the page edge, the running-head rule sits
% 112 pt below the top of the page and the last line 100 pt above the foot.
\usepackage[french]{babel}
\frenchsetup{StandardItemLabels=true}
\usepackage{amsmath,amssymb,amsthm}
\usepackage{enumitem}
\usepackage{fancyhdr}
\usepackage{tikzflora}

\pagestyle{fancy}
\fancyhf{}
\fancyhead[R]{\small 3\quad\textsc{Structures algébriques}}
\fancyfoot[R]{\small\thepage}
\renewcommand\headrulewidth{0pt}

\newtheorem{definition}{Définition}[section]
\newtheorem{proposition}[definition]{Proposition}
\theoremstyle{definition}
\newtheorem*{exemples}{Exemples}
\renewcommand\thesection{3.\arabic{section}}
\setcounter{section}{1}

\begin{document}

\section{Groupes}

Dans cette section, on introduit la notion de groupe, qui est l'une des
structures algébriques fondamentales en mathématiques. Les groupes
apparaissent naturellement dans de nombreux contextes, aussi bien en algèbre
qu'en géométrie, et permettent d'unifier une grande variété d'exemples.

\begin{definition}
Un \emph{groupe} est un ensemble $G$ muni d'une \emph{loi de composition
interne} $\ast : G\times G\to G$ telle que :
\begin{enumerate}
  \item \emph{(Associativité)} Pour tous $a,b,c\in G$, on a
    $(a\ast b)\ast c = a\ast(b\ast c)$.
  \item \emph{(Élément neutre)} Il existe un élément $e\in G$ tel que pour
    tout $a\in G$, $e\ast a = a\ast e = a$.
  \item \emph{(Inverse)} Pour tout $a\in G$, il existe un élément
    $a^{-1}\in G$ tel que $a\ast a^{-1} = a^{-1}\ast a = e$.
\end{enumerate}
\end{definition}

Lorsque la loi $\ast$ est commutative, on dit que le groupe est
\emph{abélien}. On note alors généralement la loi de manière additive :
$a+b$ au lieu de $a\ast b$.

\begin{exemples}\leavevmode
\begin{itemize}
  \item $(\mathbb{Z},+)$ est un groupe abélien.
  \item $(\mathbb{R}^\ast,\times) = (\mathbb{R}\setminus\{0\},\times)$ est un
    groupe abélien.
  \item Pour $n\ge 3$, le groupe symétrique $S_n$ des permutations de
    $\{1,\dots,n\}$ est un groupe non abélien.
  \item Le groupe linéaire inversible
    $GL_n(\mathbb{R}) = \{A\in M_n(\mathbb{R}) \mid \det(A)\neq 0\}$ est un
    groupe pour la multiplication des matrices, en général non abélien.
\end{itemize}
\end{exemples}

\subsection{Sous-groupes}

Soit $G$ un groupe. Un sous-ensemble $H\subset G$ est un \emph{sous-groupe}
de $G$ si et seulement s'il est stable par la loi de composition et par
passage aux inverses.

\begin{proposition}[Critère de sous-groupe]
Soit $H\subset G$ un sous-ensemble non vide. Alors $H$ est un sous-groupe de
$G$ si et seulement si, pour tous $a,b\in H$, on a
\[ a\ast b^{-1}\in H. \]
\end{proposition}

\begin{proof}
($\Rightarrow$) Si $H$ est un sous-groupe, alors par stabilité on a $a\in H$
et $b^{-1}\in H$, donc $a\ast b^{-1}\in H$.

($\Leftarrow$) Supposons que $H$ soit non vide et que, pour tous
$a,b\in H$, on ait $a\ast b^{-1}\in H$. Soit $a\in H$. Alors
$e = a\ast a^{-1}\in H$, donc $H$ contient l'élément neutre. De plus, pour
tout $a\in H$, on a $a^{-1} = e\ast a^{-1}\in H$. Enfin, si $a,b\in H$, alors
$a\ast b = a\ast(b^{-1})^{-1}\in H$. Ainsi $H$ est un sous-groupe de $G$.
\end{proof}

\begin{flushright}\small\itshape
Ad aliquid nutrire.\\ S. Ramanujan
\end{flushright}

\newpage

\subsection{Morphismes de groupes}

Soient $(G,\ast)$ et $(H,\cdot)$ deux groupes. Une application
$f : G\to H$ est un \emph{morphisme de groupes} si, pour tous $a,b\in G$,
\[ f(a\ast b) = f(a)\cdot f(b). \]

\begin{proposition}
Soit $f : G\to H$ un morphisme de groupes. Alors $f(e_G) = e_H$ et, pour
tout $a\in G$, $f(a^{-1}) = f(a)^{-1}$.
\end{proposition}

\begin{proof}
On a $f(e_G) = f(e_G\ast e_G) = f(e_G)\cdot f(e_G)$; en composant par
$f(e_G)^{-1}$ on obtient $f(e_G) = e_H$. Ensuite
$f(a)\cdot f(a^{-1}) = f(a\ast a^{-1}) = f(e_G) = e_H$.
\end{proof}

\begin{definition}
Le \emph{noyau} de $f$ est $\ker f = \{a\in G \mid f(a) = e_H\}$ et son
\emph{image} est $\operatorname{Im} f = \{f(a) \mid a\in G\}$.
\end{definition}

\begin{itemize}
  \item $\ker f$ est un sous-groupe de $G$ ;
  \item $\operatorname{Im} f$ est un sous-groupe de $H$ ;
  \item $f$ est injectif si et seulement si $\ker f = \{e_G\}$.
\end{itemize}

\end{document}
